\documentclass[11pt]{article}

\usepackage[preprint]{acl}
\usepackage{times}
\usepackage{latexsym}
\usepackage[T1]{fontenc}
\usepackage[utf8]{inputenc}
\usepackage{float}
\usepackage{microtype}
\usepackage{inconsolata}
\usepackage{graphicx}
\usepackage{booktabs}
\usepackage{array}
\usepackage{url}
\usepackage{amsmath}
\usepackage{placeins}


\title{Adaptive Context Selection for Token-Efficient Retrieval-Augmented Generation}

\author{Andreia Alexa \And Simon Backhaus Brudzewski \And Karlis Martins Auce \AND Joel Vazquez \And Sheraz Ahmad}


\begin{document}
\maketitle
\begin{abstract}
Retrieval-Augmented Generation (RAG) systems often pass a fixed number of retrieved documents to the generator, even when a query could be answered from less evidence. This increases prompt cost and latency, and may introduce distracting context. We present Safe Adaptive Context, a lightweight controller that selects and compresses retrieved evidence according to answer type, and expands the context only when the first answer appears weak or unsupported. All methods use the same TF-IDF top-K candidate list, which lets us isolate context selection from retrieval.We evaluate on SciFact, HotpotQA, and BioASQ with Mistral 7B and Llama-3.3-70B, comparing against fixed top-$k$ and heuristic baselines, and test robustness with cross-encoder reranking. 
Safe Adaptive Context reduces provider-reported token usage by 29--74\% relative to Fixed Top-10 while usually maintaining comparable automatic answer-quality scores, with fallback rates of 0--5\%. An exploratory ASQA experiment suggests that the same strategy can transfer to longer explanatory answers, reducing tokens by 59.1\% relative to Fixed Top-10. We complement the automatic evaluation with a structured human annotation study of 120 outputs from the three main datasets.
\end{abstract}

\section{Introduction}

Retrieval-Augmented Generation (RAG) improves language-model answers by supplying retrieved evidence at generation time (Lewis et al., 2020). In many practical systems, however, the retrieval depth is fixed: every query receives the same top-$k$ documents. This is operationally convenient but inefficient. A simple factual question may need only one sentence, while a multi-hop question may require evidence from several documents, and a longer explanatory question may require broader synthesis. Treating these cases identically wastes prompt tokens, increases latency, and can reduce answer quality when irrelevant context distracts the generator (Shi et al., 2023).

We ask a focused question: given the same initially retrieved candidate set, can a RAG system decide how much evidence to send to the generator, and in what form? This separates retrieval from context selection. Fixed and adaptive methods use the same top-10 retrieved list; they differ only in how that list is converted into prompt context. This lets us study context control directly, without conflating it with retriever improvements.

We propose Safe Adaptive Context, a lightweight task-aware controller for token-efficient RAG. The controller chooses between compact evidence selection and token-budget sentence packing depending on the expected answer type. It then applies answer-aware fallback: if the first generated answer appears empty, weakly grounded, or poorly supported by the selected context, the system expands the context and regenerates. The goal is not simply to minimize tokens, but to allocate context according to task structure and observed answer risk.

Our contributions are threefold. First, we introduce Safe Adaptive Context, a practical context-selection controller that requires no generator fine-tuning and no additional learned critic. Second, we evaluate it on SciFact, HotpotQA, and BioASQ across two model scales using provider-reported token counts, with ASQA as an exploratory long-form QA setting. Third, we analyze the quality--efficiency trade-off through fixed top-$k$ baselines, a rule-based heuristic budget controller, component diagnostics, cross-encoder reranking, and human annotation. Code and project materials are available online.

\section{Related Work}

\paragraph{Retrieval-augmented generation.}
RAG improves generation by retrieving external evidence and placing it in the model prompt (Lewis et al., 2020). Benchmarks such as BEIR (Thakur et al., 2021) evaluate retrieval across domains. Our work focuses on the post-retrieval step: once a ranked list has been retrieved, how much of it should enter the prompt, and in what form?

\paragraph{Adaptive retrieval.}
Prior work studies when to retrieve or when to retrieve more. FLARE retrieves when generation appears uncertain (Jiang et al., 2023a), while Self-RAG trains models to critique retrieval and generation through reflection tokens (Asai et al., 2023). Closest to our setting, Adaptive-$k$ selects the number of passages from the query-passage similarity distribution without fine-tuning or extra LLM calls (Taguchi et al., 2025). Our method follows the same lightweight motivation, but extends the decision from how many passages to include to how the selected evidence should be structured.

\paragraph{Context compression.}
Prompt-compression methods such as LLMLingua, Selective Context, and RECOMP reduce input length by selecting, deleting, or rewriting context (Jiang et al., 2023b; Li et al., 2023; Xu et al., 2023). Our compression method is simpler and computationally cheaper: it selects query-relevant sentence neighborhoods from retrieved documents and does not require an additional model call. A key difference is that we treat compression as task-dependent, since exact short-answer strings and multi-hop evidence can be lost through overly aggressive compression.

\paragraph{Token efficiency in RAG.}
Prior work shows that irrelevant or poorly placed context can hurt LLM performance when useful information is mixed with distractors or buried in long prompts (Shi et al., 2023; Wu et al., 2024). This motivates treating prompt context as a budgeted resource rather than always passing the full top-$k$ list. We extend this idea to a task-aware context-control setting, measuring answer quality, provider-reported token use, and fallback rate.

\section{Method}

\subsection{Shared Retriever}

The retrieval stage is intentionally lightweight because the project's primary
question is context allocation rather than retriever optimization. All methods use
the same TF-IDF top-10 candidate list, so differences come from context selection
rather than retrieval. Fixed baselines use prefixes of this list with $k \in
\{3,5,7,10\}$, while adaptive methods decide how much of the candidate set to
include.

For a query $q$, the retriever returns a ranked candidate list
$R_q = [(d_1,s_1),\ldots,(d_{10},s_{10})]$, where $d_i$ is the document at rank $i$
and $s_i$ is its TF-IDF cosine similarity score. All context-selection methods
operate on $R_q$. Retrieval quality is reported with nDCG@10 and MRR@10 using
dataset qrels; TF-IDF and cosine similarity details are provided in
Appendix~\ref{app:retriever-details}.



\subsection{Heuristic Budget Baseline}

Before presenting the main model, we describe the Heuristic Rules baseline (heuristic\_rules), which applies a set of retrieval-score-based rules to select context size:


A sharp drop after rank 3 suggests that the top few documents are much stronger than the rest. A moderate drop after rank 5 suggests medium confidence. If neither gap is large, the query is treated as ambiguous and receives more context.


The query-length-aware variant adjusts these thresholds using the number of query tokens. Short queries, defined as six tokens or fewer, use tighter thresholds because they are often underspecified. Longer claim-style queries can use slightly lower thresholds because they usually contain more specific evidence cues. The heuristic is explainable and adds essentially no latency beyond retrieval.


\subsection{Budget Prediction Component}

For adaptive evidence-style settings, Safe Adaptive Context uses a lightweight budget predictor to choose an initial document budget before generation. The predictor is implemented as a nearest-centroid classifier and is trained on a development split to imitate a minimum-sufficient oracle. The oracle selects the smallest budget in \{3, 5, 8, 10\} that reaches at least 95\% of the best available development-set utility for that query, where utility rewards answer/context quality proxies while penalizing token use.


At inference time, the predictor uses only retrieval-side features available before any LLM call, including query length, lexical diversity, top retrieval scores, score gaps, score-mass features, score ratios, entropy, and normalized entropy. This keeps the initial budget decision cheap, interpretable, and independent of the generator. The predictor is not treated as a separate final system; rather, it provides the first-pass budget used by the adaptive controller before answer-aware fallback is applied.


\subsection{Evidence Compression: n-gram Neighbor Selection}

For compact evidence modes, each selected document is split into sentences. Candidate sentences are scored using a transparent weighted function that combines lexical overlap, rare query-term matches, phrase overlap, token density, and a small position bonus. The highest-scoring sentence and its immediate neighbors are retained, and the first sentence of an abstract is prepended when useful for local topic context. This produces compact evidence spans rather than isolated sentences, preserving nearby explanatory context while reducing prompt length. The full scoring formula is provided in Appendix~\ref{app:evidence-scoring}.


\subsection{Answer-aware Fallback}

Answer-aware Fallback is the final Safe Adaptive Context policy. For evidence-style prompts, the first pass uses the predicted adaptive budget with compact evidence spans. If the generated answer appears risky, the controller progressively expands the context: first to fuller evidence at the same budget, then to a larger budget, and finally to full top-10 context. For concise QA prompts, the first pass instead uses token-budget sentence packing over the top-10 pool, with fallback to full top-10 only when the answer shape is clearly broken.


The risk checker uses only runtime signals, including empty or weak-answer phrases, answer-context overlap, context-query coverage, answer-query overlap, short-answer grounding, and verbosity. Fallback cost is counted cumulatively, including failed first-pass attempts. The checker is intended as a heuristic recovery mechanism, not as a guarantee of factual correctness.


\section{Experimental Setup}

\subsection{Datasets}

Our main evaluation uses three BEIR-style datasets with 100 held-out evaluation queries each: SciFact, HotpotQA, and BioASQ. SciFact contains biomedical claim-verification queries, HotpotQA contains multi-hop general-domain questions, and BioASQ contains biomedical factoid questions. The repository stores simplified JSONL files with documents, query text, gold relevant document identifiers, and reference answers. Because SciFact provides gold evidence labels but not QA-style reference answers, we generated synthetic SciFact references with gpt-oss-120b for answer-level metrics only; retrieval metrics still use the original qrels. We additionally report ASQA as an exploratory long-form QA experiment to test whether the same context-control strategy transfers to longer, more explanatory answers.


\subsection{Models}

Mistral 7B (referred to as Mistral throughout) is run locally via the Ollama OpenAI-compatible API at temperature 0 with a maximum of 80 completion tokens. Token counts are taken from provider-reported usage fields in the API response; runs that fail to obtain provider token counts are rejected. This ensures that comparisons are based on actual LLM billing units rather than local heuristic estimates.


Llama-3.3-70B-Instruct (referred to as Llama-70B throughout) is accessed through a hosted API at temperature 0. For the main three-dataset evaluation, Llama-70B uses the same 80-token completion limit as Mistral. The exploratory ASQA run uses a longer-answer setting appropriate for explanatory QA. The larger model provides a check that the efficiency gains observed with Mistral are not model-specific.


\subsection{Compared Methods}

We compare the proposed controller against closed-book generation, fixed-context RAG baselines, and a strong explainable heuristic budget controller. Fixed Top-10 is treated as the expensive fixed-context baseline, not as a guaranteed quality upper bound.


We also report three diagnostic extensions: component ablations, cross-encoder reranking, and the exploratory ASQA long-form run. These tests examine which parts of the controller matter, whether the method still helps under stronger ranking, and whether it transfers to longer explanatory answers. They are reported separately from the main three-dataset comparison.


\subsection{Evaluation Metrics}

We evaluate answer quality with token F1, answer coverage, and semantic similarity; retrieval quality with nDCG@10 and MRR@10; and efficiency with provider-reported total tokens, token reduction relative to Fixed Top-10, and fallback rate. For ASQA, we interpret token F1 together with coverage and semantic similarity because longer explanatory answers are less reliably captured by lexical overlap alone.

\section{Results}

\subsection{Main Results}

To keep the main text focused, Table~\ref{tab:table_1_safe_adaptive_context_compared_w} reports the central comparison between Safe Adaptive Context and the Fixed Top-10 baseline across all dataset-model pairs. The full method ladder, including No Retrieval, Fixed Top-3, Fixed Top-5, Fixed Top-7, Fixed Top-10, Heuristic Rules, and Safe Adaptive Context, is reported by dataset in Appendix~\ref{app:full-main-results}.


\begin{table*}[t]
\centering
\small
\caption{Safe Adaptive Context compared with fixed top-10 full-document RAG.}
\label{tab:table_1_safe_adaptive_context_compared_w}
\resizebox{\linewidth}{!}{%
\begin{tabular}{lllrrrrrr}
\toprule
Dataset & Model & Fixed F1 & Safe F1 & Fixed Semantic & Safe Semantic & Fixed Tokens & Safe Tokens & Token vs Fixed \\
\midrule
SciFact & Mistral & 0,204 & 0,253 & 0,364 & 0,426 & 3.781 & 985 & -73,9\% \\
SciFact & Llama-70B & 0,262 & 0,265 & 0,569 & 0,569 & 3.412 & 933 & -72,6\% \\
HotpotQA & Mistral & 0,516 & 0,540 & 0,529 & 0,554 & 1.762 & 1.240 & -29,6\% \\
HotpotQA & Llama-70B & 0,765 & 0,737 & 0,768 & 0,740 & 1.537 & 1.066 & -30,6\% \\
BioASQ & Mistral & 0,257 & 0,339 & 0,606 & 0,758 & 3.889 & 2.011 & -48,3\% \\
BioASQ & Llama-70B & 0,344 & 0,342 & 0,774 & 0,779 & 3.496 & 1.668 & -52,3\% \\
\bottomrule
\end{tabular}%
}
\end{table*}

Token counts are provider-reported average total tokens per query for the final reported runs. Table~\ref{tab:table_2_safe_adaptive_context_compared_w} compares Safe Adaptive Context with the Heuristic Rules baseline, testing whether the systematic adaptive controller adds value beyond a simple score-gap budget policy.

Safe Adaptive Context reduces token use in every setting. The largest reductions occur on SciFact, where evidence is often concentrated in a small number of biomedical sentences. HotpotQA yields smaller but still substantial savings because multi-hop questions require broader coverage. BioASQ falls between these cases.

In terms of automatic answer quality, Safe Adaptive Context generally remains close to Fixed Top-10 and sometimes improves over it. However, several differences are small and should not be interpreted as robust superiority without statistical testing. The main result is therefore best understood as a favorable efficiency--quality trade-off: Safe Adaptive often achieves comparable automatic scores while using substantially fewer provider-reported tokens.

The No Retrieval baseline also shows that token F1 alone can be misleading. On SciFact with Mistral, No Retrieval obtains higher token F1 than Fixed Top-10 and Safe Adaptive Context, even though its semantic similarity and answer coverage are lower. This suggests that SciFact token F1 is sensitive to answer phrasing and label-like output formats, and may reward short generic or claim-style answers even when retrieved evidence improves grounding. For this reason, we interpret token F1 together with semantic similarity, coverage, retrieval metrics, and human annotation rather than treating it as a complete measure of factual correctness.

A similar caution applies to BioASQ, where semantic similarity is relatively high even for No Retrieval. Biomedical terminology and generic domain-relevant answers may receive high embedding similarity despite being incomplete or insufficiently specific. This motivates the human annotation study in Section~\ref{sec:human-annotation}, which provides a complementary check on whether automatic scores correspond to factual answer quality.

\begin{table*}[t]
\centering
\small
\caption{Safe Adaptive Context compared with the Heuristic Rules baseline.}
\label{tab:table_2_safe_adaptive_context_compared_w}
\resizebox{\linewidth}{!}{%
\begin{tabular}{lllrrrrrr}
\toprule
Dataset & Model & Heuristic F1 & Safe F1 & Heuristic Semantic & Safe Semantic & Heuristic Tokens & Safe Tokens & Token vs Heuristic \\
\midrule
SciFact & Mistral & 0,265 & 0,253 & 0,437 & 0,426 & 2.652 & 985 & -62,9\% \\
SciFact & Llama-70B & 0,266 & 0,265 & 0,575 & 0,569 & 2.234 & 933 & -58,2\% \\
HotpotQA & Mistral & 0,482 & 0,540 & 0,497 & 0,554 & 1.073 & 1.240 & 15,5\% \\
HotpotQA & Llama-70B & 0,708 & 0,737 & 0,711 & 0,740 & 945 & 1.066 & 12,8\% \\
BioASQ & Mistral & 0,333 & 0,339 & 0,745 & 0,758 & 2.560 & 2.011 & -21,4\% \\
BioASQ & Llama-70B & 0,349 & 0,342 & 0,780 & 0,779 & 2.126 & 1.668 & -21,5\% \\
\bottomrule
\end{tabular}%
}
\end{table*}


\subsection{Ablation and compression diagnostic}

The main comparison shows the overall quality--efficiency trade-off, but it does not by itself explain which parts of the controller are responsible for the gains. We therefore report two diagnostic analyses. The first is a component ablation on SciFact with Llama-70B over 25 held-out queries, which isolates simplified versions of the Safe Adaptive pipeline. The second compares full context, compressed context, heuristic budgeting, and Safe Adaptive on the final Llama-70B runs for SciFact. The full compression diagnostic across SciFact and HotpotQA is provided in Appendix~\ref{app:diagnostics}. In the main text, we focus on SciFact/Llama-70B as a controlled diagnostic setting because it isolates the effect of compression and budget selection without introducing a second task structure.


\begin{table*}[t]
\centering
\small
\caption{Component ablation on SciFact with Llama-70B over 25 held-out queries.}
\label{tab:table_3_component_ablation_on_scifact_wi}
\resizebox{\linewidth}{!}{%
\begin{tabular}{lllrrrr}
\toprule
Variant & nDCG@10 & F1 & Semantic & Tokens & Token Reduction & Fallback \\
\midrule
Fixed Top-10 & 0,558 & 0,247 & 0,489 & 3.325 & 0,0\% & 0,0\% \\
Adaptive Compact Only & 0,484 & 0,248 & 0,481 & 736 & 77,9\% & 0,0\% \\
Adaptive Full Only & 0,484 & 0,240 & 0,485 & 1.097 & 67,0\% & 0,0\% \\
Fixed Top-5 Compact + Fallback & 0,532 & 0,258 & 0,489 & 1.209 & 63,6\% & 4,0\% \\
Full Safe Adaptive Context & 0,484 & 0,247 & 0,477 & 780 & 76,5\% & 4,0\% \\
\bottomrule
\end{tabular}%
}
\end{table*}

The component ablation shows that compact evidence is not simply a lossy shortcut. Adaptive Compact Only matches Fixed Top-10 F1 while using 77.9\% fewer tokens, whereas Adaptive Full Only is less efficient and slightly lower in F1. This suggests that compression can remove distracting context rather than merely shortening the prompt. The full Safe Adaptive pipeline preserves Fixed Top-10 F1 with a 76.5\% token reduction, while fallback is triggered on only 4\% of queries.


\begin{table}[t]
\centering
\small
\caption{Compression and budget diagnostic on Llama-70B.}
\label{tab:table_4_compression_and_budget_diagnosti}
\resizebox{\linewidth}{!}{%
\begin{tabular}{lrrrr}
\toprule
Method & F1 & Semantic & Tokens & Token Reduction \\
\midrule
Fixed Full Context & 0,262 & 0,569 & 3.412 & 0,0\% \\
Compressed Fixed Full Context & 0,257 & 0,571 & 2.150 & 37,0\% \\
Heuristic Rules & 0,266 & 0,575 & 2.234 & 34,5\% \\
Safe Adaptive Context & 0,265 & 0,569 & 933 & 72,6\% \\
\bottomrule
\end{tabular}%
}
\end{table}

The compression diagnostic separates three effects on the same final SciFact/Llama-70B run: using the full Fixed Top-10 context, compressing that same full context, and applying an adaptive budget policy. Compressed full context preserves almost all F1 while reducing tokens by 37.0\%, showing that much of the full prompt can be removed without damaging answer quality. Heuristic Rules achieves slightly higher F1 than Fixed Top-10 with a 34.5\% token reduction, but Safe Adaptive reaches nearly the same F1 while reducing tokens by 72.6\%. This suggests that the main gain comes from combining compact evidence selection with adaptive budget control, rather than from compression alone.


\subsection{Cross-Encoder Reranking}

The main experiments use TF-IDF retrieval so that the context-control policy can be tested in a simple and transparent setting. To check whether the gains depend on this particular ranking method, we also rerank the same top-10 candidate pool with a cross-encoder before applying the fixed and adaptive context policies. This experiment is a robustness check rather than a new retrieval benchmark: the question is whether Safe Adaptive Context still preserves answer quality and token savings when the ordering of evidence changes.


\begin{table*}[t]
\centering
\small
\caption{Robustness of Safe Adaptive Context under TF-IDF and cross-encoder reranking. Fixed F1 is the Fixed Top-10 score under the same retriever. Token reduction is measured against Fixed Top-10 for that retriever.}
\label{tab:cross_encoder}

\begin{tabular}{llrrrrrr}
\toprule
Dataset & Retriever & Fixed F1 & Safe F1 & Safe F1 vs Fixed & Safe Tokens & Token Reduction & Safe nDCG \\
\midrule
SciFact & TF-IDF & 0.262 & 0.265 & +0.003 & 933 & 72.6\% & 0.560 \\
SciFact & Cross-enc. & 0.261 & 0.265 & +0.004 & 939 & 72.5\% & 0.624 \\
HotpotQA & TF-IDF & 0.765 & 0.737 & -0.028 & 1,066 & 30.6\% & 0.742 \\
HotpotQA & Cross-enc. & 0.764 & 0.751 & -0.013 & 1,109 & 26.1\% & 0.829 \\
BioASQ & TF-IDF & 0.344 & 0.342 & -0.002 & 1,668 & 52.3\% & 0.825 \\
BioASQ & Cross-enc. & 0.338 & 0.342 & +0.004 & 814 & 76.7\% & 0.672 \\
\bottomrule
\end{tabular}%
\end{table*}

Table~\ref{tab:cross_encoder} summarizes the cross-encoder robustness check. Overall, Safe Adaptive Context is not tied to TF-IDF ordering. Cross-encoder reranking changes the operating point, but the quality-efficiency trade-off remains favorable across all three datasets. Full TF-IDF and cross-encoder method-level results are provided in Appendix Table~\ref{tab:appendix_table_a5_full_tf_idf_and_cross_encod}.


\subsection{Exploratory Long-Form Evaluation on ASQA}

The main evaluation focuses on compact claim-verification, factoid, and multi-hop QA settings. To test whether the same context-control strategy transfers to longer explanatory answers, we also run an exploratory ASQA experiment with Llama-70B. ASQA is treated separately from the main three-dataset benchmark because answer quality is less well captured by lexical overlap alone. We therefore interpret token F1 together with answer coverage and semantic similarity.


\begin{table*}[t]
\centering
\small
\setlength{\tabcolsep}{4pt}
\caption{Exploratory ASQA long-form evaluation with Llama-70B. Token reduction is measured relative to Fixed Top-10.}
\label{tab:table_6_exploratory_asqa_long_form_evalu}
\resizebox{0.82\linewidth}{!}{%
\begin{tabular}{lllrrrr}
\toprule
Method & F1 & Coverage & Semantic & Tokens & Reduction & Fallback \\
\midrule
No Retrieval & 0,203 & 0,159 & 0,387 & 113 & 90,8\% & 0,0\% \\
Fixed Top-3 & 0,398 & 0,395 & 0,772 & 422 & 65,7\% & 0,0\% \\
Fixed Top-5 & 0,405 & 0,417 & 0,783 & 625 & 49,2\% & 0,0\% \\
Fixed Top-7 & 0,407 & 0,423 & 0,790 & 854 & 30,6\% & 0,0\% \\
Fixed Top-10 & 0,411 & 0,419 & 0,786 & 1.229 & 0,0\% & 0,0\% \\
Heuristic Rules & 0,393 & 0,414 & 0,782 & 629 & 48,8\% & 0,0\% \\
Safe Adaptive Context & 0,417 & 0,411 & 0,782 & 503 & 59,1\% & 6,0\% \\
\bottomrule
\end{tabular}%
}
\end{table*}

Table~\ref{tab:table_6_exploratory_asqa_long_form_evalu} reports the ASQA long-form stress test. Safe Adaptive achieves the highest F1, outperforming the Heuristic Rules baseline while reducing tokens by 59.1\% relative to Fixed Top-10. We treat ASQA as exploratory because long-form QA is harder to judge with lexical overlap alone.


\subsection{Summary Across Datasets}

Across datasets, Safe Adaptive Context reduces tokens substantially while retaining at least comparable answer quality on average; a compact dataset-level summary is provided in Table~\ref{tab:table_7_dataset_level_summary_of_safe_ad}. The variation supports the main design claim: context budgets should depend on task structure. Concentrated evidence tasks such as SciFact permit aggressive compression, while multi-hop tasks such as HotpotQA require wider evidence coverage.


\section{Human Annotation Study}
\label{sec:human-annotation}

Automatic metrics such as Token F1 and semantic similarity are useful for large-scale comparison, but they do not fully measure factual correctness. Token F1 can penalize correct paraphrases, while embedding similarity can remain high for answers that are related but still incomplete or wrong. We therefore include a structured human annotation study to check whether the automatic trends align with human judgments.

\subsection{Annotation Protocol}

Each annotator sees the query, reference answer, model answer, and retrieved evidence. Ratings are based on answer content rather than wording style. We use four labels: CORRECT, PARTIALLY\_CORRECT, INCORRECT, and NOT\_ENOUGH\_INFO; full definitions are provided in Appendix~\ref{app:annotation-protocol}.

\subsection{Sample Construction}

The annotation set contains 120 Safe Adaptive outputs, with 40 examples each from SciFact, HotpotQA, and BioASQ. SciFact and BioASQ are stratified across Mistral and Llama-70B outputs, while HotpotQA covers the observed output distribution. Sampling is stratified by semantic similarity, so the set includes low-, middle-, and high-scoring examples. Two annotators independently label each item. Full sample-construction details are provided in Appendix~\ref{app:annotation-protocol}.


\subsection{Annotation Interface and Agreement}

The browser-based workform is self-contained and requires no server. It supports local saving, CSV export/import, filtering, search, and dataset tabs. Annotators work independently. Agreement is computed with a Cohen's kappa calculator that matches items by ID and lists disagreements for adjudication.

Across the 120 matched items, annotators agreed on 84 examples and disagreed on 36, yielding 70.0\% raw agreement and Cohen's $\kappa = 0.477$. This indicates moderate agreement. The value is acceptable for factual RAG evaluation, but it also shows that the task contains substantial borderline cases, especially between CORRECT, PARTIALLY\_CORRECT, and INCORRECT labels. We therefore interpret the human-label results as a complementary quality signal rather than as an absolute ground truth.

\subsection{Annotation Results}

Table~\ref{tab:annotation_all_labels} reports the label distribution across the 120 annotated examples, counting both annotators' labels. Overall, annotators assigned 147 CORRECT labels (61.3\%), 34 PARTIALLY\_CORRECT labels (14.2\%), 45 INCORRECT labels (18.8\%), and 14 NOT\_ENOUGH\_INFO labels (5.8\%).

\begin{table}[t]
\centering
\small
\caption{Distribution of individual annotation labels across both annotators.}
\label{tab:annotation_all_labels}
\resizebox{\linewidth}{!}{%
\begin{tabular}{lrrrr}
\toprule
Dataset & Correct & Partial & Incorrect & Not Enough Info \\
\midrule
SciFact & 49 & 15 & 14 & 2 \\
HotpotQA & 54 & 1 & 17 & 8 \\
BioASQ & 44 & 18 & 14 & 4 \\
Overall & 147 & 34 & 45 & 14 \\
\bottomrule
\end{tabular}%
}
\end{table}

For metric validation, we use the 84 agreed examples as a conservative subset with unambiguous human labels. Table~\ref{tab:annotation_agreed_by_dataset} shows that 65 of 84 agreed examples were labeled CORRECT, and 71 of 84 were labeled either CORRECT or PARTIALLY\_CORRECT. This subset is useful for checking whether automatic scores align with clear human judgments, but it is biased toward easier cases. The 36 disagreed examples are likely to contain many of the most informative borderline cases, where answers are partially grounded, incomplete, or phrased differently from the reference.

\begin{table}[t]
\centering
\small
\caption{Human labels for agreed annotation examples only.}
\label{tab:annotation_agreed_by_dataset}
\resizebox{\linewidth}{!}{%
\begin{tabular}{lrrrr}
\toprule
Dataset & Agreed Items & Correct & Partial & Incorrect \\
\midrule
SciFact & 27 & 20 & 2 & 5 \\
HotpotQA & 30 & 26 & 0 & 4 \\
BioASQ & 27 & 19 & 4 & 4 \\
Overall & 84 & 65 & 6 & 13 \\
\bottomrule
\end{tabular}%
}
\end{table}

Table~\ref{tab:annotation_metric_alignment} compares agreed human labels against automatic metric buckets. Human correctness increases strongly with both Token F1 and semantic similarity. All high-F1 agreed examples were judged correct, and 90.9\% of medium-F1 examples were judged correct. Semantic similarity shows a similar pattern: 95.9\% of high-semantic examples were judged correct.

\begin{table*}[t]
\centering
\small
\setlength{\tabcolsep}{4pt}
\caption{Alignment between automatic metric buckets and agreed human labels. Percentages are computed within each bucket.}
\label{tab:annotation_metric_alignment}
\resizebox{0.78\linewidth}{!}{%
\begin{tabular}{llrrrr}
\toprule
Metric & Bucket & Items & Correct & Partially Correct & Incorrect \\
\midrule
Token F1 & Low (0.00--0.33) & 33 & 48.5\% & 18.2\% & 33.3\% \\
Token F1 & Medium (0.34--0.66) & 22 & 90.9\% & 0.0\% & 9.1\% \\
Token F1 & High (0.67--1.00) & 29 & 100.0\% & 0.0\% & 0.0\% \\
\midrule
Semantic & Low (0.00--0.49) & 23 & 47.8\% & 13.0\% & 39.1\% \\
Semantic & Medium (0.50--0.74) & 12 & 58.3\% & 25.0\% & 16.7\% \\
Semantic & High (0.75--1.00) & 49 & 95.9\% & 0.0\% & 4.1\% \\
\bottomrule
\end{tabular}%
}
\end{table*}

Correct answers also have substantially higher average automatic scores than partially correct or incorrect answers. Among agreed examples, CORRECT answers average 0.615 F1, 0.775 semantic similarity, and 0.862 answer coverage. PARTIALLY\_CORRECT answers average 0.190 F1, 0.473 semantic similarity, and 0.309 coverage, while INCORRECT answers average 0.143 F1, 0.291 semantic similarity, and 0.228 coverage.

\subsection{Annotation Outcomes}

The annotation results support the automatic evaluation while clarifying its limits. High Token F1 and high semantic similarity are strongly associated with human-correct answers, which suggests that the automatic metrics are useful directional indicators of answer quality. However, low-F1 examples are mixed: 48.5\% are still judged correct and 18.2\% partially correct. This confirms that lexical overlap can underestimate acceptable answers when they are paraphrased or phrased differently from the reference.

The dataset-level patterns are also informative. HotpotQA has the highest proportion of agreed CORRECT labels, with 26 of 30 agreed examples marked correct, consistent with strong automatic performance and high F1 retention. BioASQ has the largest number of PARTIALLY\_CORRECT agreed labels, matching the expectation that biomedical answers often contain relevant information but may be incomplete or broader than the concise reference. SciFact receives few NOT\_ENOUGH\_INFO labels in the full annotation distribution, suggesting that annotators usually judged its outputs as supported, contradicted, or incomplete rather than impossible to assess.

Overall, the human study validates our use of automatic metrics as useful but incomplete evaluation signals. The results support reporting Token F1 together with semantic similarity, answer coverage, and human annotation rather than relying on any single metric.

\section{Discussion}

The central finding is that context size is a decision variable rather than a fixed setting. Fixed Top-10 is simple, but not uniformly best: smaller or compressed contexts often improve Mistral results on SciFact and BioASQ, suggesting that weaker models are more sensitive to distracting passages. For Llama-70B, the main benefit is cost rather than quality.

The gains follow task structure. SciFact is easiest to compress because many claims can be checked from a small number of evidence sentences. BioASQ is intermediate, since biomedical abstracts are relevant but broad. HotpotQA is the limiting case: multi-hop questions often require evidence from multiple documents, so Safe Adaptive saves fewer tokens and relies on broader sentence packing.

The results also suggest that compression should be controlled rather than applied uniformly. Compact evidence works well when relevant information is concentrated, as in many SciFact examples, but the HotpotQA results show that multi-hop questions often need broader coverage. This supports the central design of Safe Adaptive Context: the system does not simply compress every prompt, but chooses between compact evidence, sentence packing, and fallback depending on the task and answer-risk signals.

The method depends on retrieval quality. If relevant evidence appears late in the top-10 list, aggressive budgeting can remove it. Cross-encoder reranking partly mitigates this by moving useful evidence earlier, while preserving the same quality--efficiency pattern. Better retrieval therefore complements adaptive budgeting rather than replacing it.

The ASQA result suggests that context control can transfer to longer answers, but it should not be read as evidence that long-form RAG can always be aggressively compressed. Coverage is slightly lower than Fixed Top-10, and long-form answer quality is difficult to capture with lexical metrics alone.

Overall, Safe Adaptive Context is best viewed as a lightweight operating policy: it spends tokens when the task or ranking suggests uncertainty, and saves them when evidence is concentrated. Its thresholds are not universal, and the human annotation results show why automatic metrics should be interpreted cautiously.



\section{Conclusion}

Safe Adaptive Context reduces RAG token use by adapting context size and expanding only on risky answers. Across three main datasets and two model scales, it saves 29.6--73.9\% of tokens relative to Fixed Top-10 while usually achieving comparable automatic answer-quality scores. In the one setting with the largest observed decrease, HotpotQA with Llama-70B, Safe Adaptive trades a modest F1 drop for a 30.6\% token reduction. Ablations and cross-encoder reranking suggest that the gains come from combining adaptive budgeting with lexical evidence compression, rather than from TF-IDF ordering alone. An exploratory ASQA run suggests that the approach may transfer to longer explanatory answers, where Safe Adaptive improves token F1 while reducing token use by 59.1\%.

The main limitation is that automatic metrics miss some aspects of factual correctness, especially for biomedical and long-form answers. Token F1 can reward short answer forms or penalize correct paraphrases, while semantic similarity may over-credit domain-related but incomplete answers. Human annotation helps identify these cases, but the moderate inter-annotator agreement also shows that factual RAG evaluation contains genuine ambiguity. Future work should evaluate stronger retrievers, report uncertainty estimates, calibrate fallback thresholds, and include adjudicated human labels for borderline cases.

\section*{References}

Asai, A., Wu, Z., Wang, Y., Sil, A., \& Hajishirzi, H. (2023). Self-RAG: Learning to retrieve, generate, and critique through self-reflection. arXiv preprint arXiv:2310.11511.


Jiang, Z., Xu, F. F., Gao, L., Sun, Z., Liu, Q., Dwivedi-Yu, J., Yang, Y., Callan, J., \& Neubig, G. (2023a). Active retrieval augmented generation. arXiv preprint arXiv:2305.06983.


Jiang, H., Wu, Q., Lin, C.-Y., Yang, P., \& Qiu, X. (2023b). LLMLingua: Compressing prompts for accelerated inference of large language models. arXiv preprint arXiv:2310.05736.


Li, Y., Dong, B., Lin, C., \& Guerin, F. (2023). Compressing context to enhance inference efficiency of large language models. arXiv preprint arXiv:2310.06201.


Shi, F., Chen, X., Misra, K., Scales, N., Dohan, D., Chi, E., Schärli, N., \& Zhou, D. (2023). Large language models can be easily distracted by irrelevant context. Proceedings of ICML 2023.


Thakur, N., Reimers, N., Rücklé, A., Srivastava, A., Gurevych, I. (2021). BEIR: A heterogeneous benchmark for zero-shot evaluation of information retrieval models. arXiv preprint arXiv:2104.08663.


Taguchi, Y., Maekawa, S., \& Bhutani, N. (2025). Efficient Context Selection for Long-Context QA: No Tuning, No Iteration, Just Adaptive-k. arXiv preprint arXiv:2506.08479.


Xu, F. F., Shi, W., \& Choi, E. (2023). RECOMP: Improving retrieval-augmented LMs with context compression and selective augmentation. arXiv preprint arXiv:2310.04408.


Wu, W., Wang, Y., Xiao, G., Peng, H., \& Fu, Y. (2024). Retrieval head mechanistically explains long-context factuality. arXiv preprint arXiv:2404.15574.


\appendix
\section{Project Statements}
\subsection*{Author Contributions}

  \begin{center}
  \small
  \begin{tabular}{p{0.22\linewidth}p{0.70\linewidth}}
  \toprule
  Author & Contributions \\
  \midrule
  Joel Vazquez &
  Project lead. Proposed, developed and tested the main adaptive-context idea; 
  tested evaluation metrics; coordinated implementation
  direction around the other methods and dataset selection; 
  led report writing; 
  updated and coordinated the final presentation. \\
  \midrule
  Andreia Alexa &
  Contributed to the nDCG@10 and MRR@10 metrics, fixed Top-$k$ baselines; 
  Heuristic Rules method design and testing; 
  Annotation UI web-form design; annotation data extraction and guidelines creation; 
  report review and editing; HTML presentation design and maintenance; 
  slides on setup, Top-$k$, evaluation, heuristic methods, and results; \\
  \midrule
  Simon Backhaus Brudzewski &
  Annotated samples for the human evaluation study; implemented the semantic similarity metric; calculated
Cohen's kappa; interpreted annotation results; contributed the
human annotation analysis; prepared annotation slides and reviewed the HTML presentation.  \\
  \midrule
  Karlis 
  
  Martins Auce &
  Developed and tested the cross-encoder re-ranking; ran related experiments; 
  contributed to the presentation. \\
  \midrule
  Sheraz 
  
  Ahmad &
  Annotated samples for the human evaluation study. Dataset preparation for the main experiments; contributed to the presentation.
  \\
  \bottomrule
  \end{tabular}
  \end{center}

\subsection*{Ethics Statement}
The experiments use benchmark-style datasets and generated answers. Because SciFact and BioASQ are biomedical, outputs should not be treated as medical advice or clinical evidence. We evaluate the system as an information retrieval and question-answering method, not as a clinical decision tool.


\subsection*{Reproducibility}
Code and project materials are available online.\footnote{\url{https://github.com/Joel-Vazquez-Lopez/adaptive-rag-token-efficiency}} The implementation is under \texttt{src/adaptive\_retrieval/}; experiment runners are under \texttt{scripts/}; saved final tables and ablations are under \texttt{saved\_results/}; and annotation materials are under \texttt{annotation\_workform/}. Final reported runs used temperature 0 and required provider-reported token counts.

\subsection*{AI Contributions}
We used AI writing and coding assistants to support drafting, editing, LaTeX formatting, table organization, presentation preparation, and parts of the coding structure. The authors made all experimental design decisions, reviewed the code and system behavior to ensure that the implemented pipeline matched the intended methods, checked all results and interpretations, and edited the final wording and reported claims.

\section*{Limitations}
The evaluation is intentionally controlled rather than exhaustive. We use 100 held-out queries per main dataset and a simple TF-IDF retriever to isolate context selection from retriever optimization. Stronger retrievers may change the operating point. In addition, SciFact answer-level metrics use synthetic gpt-oss-120b reference answers because the dataset does not provide QA-style natural-language answers, so we interpret these metrics together with retrieval scores and human annotation.


\clearpage
\onecolumn
\section{Implementation Details and Full Main Results}

\subsection{Retriever Details}
\label{app:retriever-details}

For a corpus of $N$ documents, inverse document frequency is computed as:

\begin{equation}
\mathrm{idf}(t) = \log\left(\frac{1 + N}{1 + \mathrm{df}(t)}\right) + 1,
\end{equation}

where $\mathrm{df}(t)$ is the number of documents containing term $t$. Query and document vectors are computed in the same TF-IDF space, and documents are ranked by cosine similarity:

\begin{equation}
\mathrm{sim}(q,d) = \frac{q \cdot d}{\lVert q \rVert \lVert d \rVert}.
\end{equation}

\subsection{Evidence Compression Scoring}
\label{app:evidence-scoring}

For compact evidence modes, candidate sentences are scored as:

\begin{equation}
\begin{split}
\mathrm{score}(q,s,i) ={}& 0.40 \cdot \mathrm{overlap}(q,s)
+ 0.25 \cdot \mathrm{rare}(q,s) \\
&+ 0.20 \cdot \mathrm{phrase}(q,s)
+ 0.10 \cdot \mathrm{density}(q,s) \\
&+ 0.05 \cdot (1 / (1 + i)).
\end{split}
\end{equation}

Here, $q$ is the query, $s$ is a candidate sentence, and $i$ is the sentence position in the document. The overlap term captures unigram query-term matches, rare gives higher weight to less frequent query terms, phrase captures bigram and trigram matches, density rewards compact sentences with many query-relevant tokens, and the position bonus mildly favors earlier sentences.

\subsection{Full Main Results by Dataset}
\label{app:full-main-results}
\subsubsection{SciFact}
Table~\ref{tab:appendix_table_a1_full_scifact_results_across} reports the full SciFact method ladder for both evaluated model families, including No Retrieval, fixed-k baselines, Heuristic Rules, and Safe Adaptive Context.
\begin{table}[htbp]
\centering
\small
\caption{Full SciFact results across retrieval methods and models.}
\label{tab:appendix_table_a1_full_scifact_results_across}
\resizebox{\linewidth}{!}{%
\begin{tabular}{lllrrrrrrrr}
\toprule
Dataset & Model & Method & NDCG@10 & MRR@10 & Answer F1 & Coverage & Semantic & Tokens & Token reduction vs Top-10 & Fallback \\
\midrule
SciFact & Mistral & No Retrieval & 0,000 & 0,000 & 0,304 & 0,638 & 0,451 & 142 & 96,3\% & 0,0\% \\
SciFact & Mistral & Fixed Top-3 & 0,560 & 0,542 & 0,268 & 0,659 & 0,439 & 1.323 & 65,0\% & 0,0\% \\
SciFact & Mistral & Fixed Top-5 & 0,587 & 0,557 & 0,265 & 0,676 & 0,443 & 2.106 & 44,3\% & 0,0\% \\
SciFact & Mistral & Fixed Top-7 & 0,597 & 0,560 & 0,262 & 0,674 & 0,439 & 2.898 & 23,4\% & 0,0\% \\
SciFact & Mistral & Fixed Top-10 & 0,602 & 0,563 & 0,204 & 0,545 & 0,364 & 3.781 & 0,0\% & 0,0\% \\
SciFact & Mistral & Heuristic Rules & 0,587 & 0,557 & 0,265 & 0,669 & 0,437 & 2.652 & 29,9\% & 0,0\% \\
SciFact & Mistral & Safe Adaptive Context & 0,560 & 0,542 & 0,253 & 0,647 & 0,426 & 985 & 73,9\% & 5,0\% \\
SciFact & Llama-70B & No Retrieval & 0,000 & 0,000 & 0,285 & 0,488 & 0,378 & 119 & 96,5\% & 0,0\% \\
SciFact & Llama-70B & Fixed Top-3 & 0,560 & 0,542 & 0,259 & 0,798 & 0,558 & 1.131 & 66,9\% & 0,0\% \\
SciFact & Llama-70B & Fixed Top-5 & 0,587 & 0,557 & 0,264 & 0,815 & 0,570 & 1.774 & 48,0\% & 0,0\% \\
SciFact & Llama-70B & Fixed Top-7 & 0,597 & 0,560 & 0,261 & 0,802 & 0,573 & 2.446 & 28,3\% & 0,0\% \\
SciFact & Llama-70B & Fixed Top-10 & 0,602 & 0,563 & 0,262 & 0,808 & 0,569 & 3.412 & 0,0\% & 0,0\% \\
SciFact & Llama-70B & Heuristic Rules & 0,587 & 0,557 & 0,266 & 0,807 & 0,575 & 2.234 & 34,5\% & 0,0\% \\
SciFact & Llama-70B & Safe Adaptive Context & 0,560 & 0,542 & 0,265 & 0,799 & 0,569 & 933 & 72,6\% & 5,0\% \\
\bottomrule
\end{tabular}%
}
\end{table}
\FloatBarrier


\subsubsection{HotpotQA}
Table~\ref{tab:appendix_table_a2_full_hotpotqa_results_acros} reports the full HotpotQA method ladder. This dataset is the multi-hop setting, where wider context remains more useful than in SciFact.
\begin{table}[htbp]
\centering
\small
\caption{Full HotpotQA results across retrieval methods and models.}
\label{tab:appendix_table_a2_full_hotpotqa_results_acros}
\resizebox{\linewidth}{!}{%
\begin{tabular}{lllrrrrrrrr}
\toprule
Dataset & Model & Method & NDCG@10 & MRR@10 & Answer F1 & Coverage & Semantic & Tokens & Token reduction vs Top-10 & Fallback \\
\midrule
HotpotQA & Mistral & No Retrieval & 0,000 & 0,000 & 0,292 & 0,347 & 0,307 & 149 & 91,5\% & 0,0\% \\
HotpotQA & Mistral & Fixed Top-3 & 0,614 & 0,790 & 0,410 & 0,477 & 0,424 & 545 & 69,1\% & 0,0\% \\
HotpotQA & Mistral & Fixed Top-5 & 0,667 & 0,806 & 0,467 & 0,508 & 0,475 & 860 & 51,2\% & 0,0\% \\
HotpotQA & Mistral & Fixed Top-7 & 0,707 & 0,812 & 0,507 & 0,570 & 0,519 & 1.192 & 32,3\% & 0,0\% \\
HotpotQA & Mistral & Fixed Top-10 & 0,742 & 0,812 & 0,516 & 0,576 & 0,529 & 1.762 & 0,0\% & 0,0\% \\
HotpotQA & Mistral & Heuristic Rules & 0,691 & 0,809 & 0,482 & 0,563 & 0,497 & 1.073 & 39,1\% & 0,0\% \\
HotpotQA & Mistral & Safe Adaptive Context & 0,742 & 0,812 & 0,540 & 0,603 & 0,554 & 1.240 & 29,6\% & 0,0\% \\
HotpotQA & Llama-70B & No Retrieval & 0,000 & 0,000 & 0,425 & 0,451 & 0,431 & 152 & 90,1\% & 0,0\% \\
HotpotQA & Llama-70B & Fixed Top-3 & 0,614 & 0,790 & 0,626 & 0,620 & 0,628 & 492 & 68,0\% & 0,0\% \\
HotpotQA & Llama-70B & Fixed Top-5 & 0,667 & 0,806 & 0,693 & 0,688 & 0,698 & 762 & 50,4\% & 0,0\% \\
HotpotQA & Llama-70B & Fixed Top-7 & 0,707 & 0,812 & 0,730 & 0,726 & 0,734 & 1.047 & 31,9\% & 0,0\% \\
HotpotQA & Llama-70B & Fixed Top-10 & 0,742 & 0,812 & 0,765 & 0,752 & 0,768 & 1.537 & 0,0\% & 0,0\% \\
HotpotQA & Llama-70B & Heuristic Rules & 0,691 & 0,809 & 0,708 & 0,696 & 0,711 & 945 & 38,5\% & 0,0\% \\
HotpotQA & Llama-70B & Safe Adaptive Context & 0,742 & 0,812 & 0,737 & 0,732 & 0,740 & 1.066 & 30,6\% & 0,0\% \\
\bottomrule
\end{tabular}%
}
\end{table}
\FloatBarrier

\subsubsection{BioASQ}
Table~\ref{tab:appendix_table_a3_full_bioasq_results_across_} reports the full BioASQ method ladder. BioASQ shows intermediate token savings between SciFact and HotpotQA.
\begin{table}[htbp]
\centering
\small
\caption{Full BioASQ results across retrieval methods and models.}
\label{tab:appendix_table_a3_full_bioasq_results_across_}
\resizebox{\linewidth}{!}{%
\begin{tabular}{lllrrrrrrrr}
\toprule
Dataset & Model & Method & NDCG@10 & MRR@10 & Answer F1 & Coverage & Semantic & Tokens & Token reduction vs Top-10 & Fallback \\
\midrule
BioASQ & Mistral & No Retrieval & 0,000 & 0,000 & 0,220 & 0,267 & 0,653 & 119 & 96,9\% & 0,0\% \\
BioASQ & Mistral & Fixed Top-3 & 0,637 & 0,940 & 0,333 & 0,436 & 0,751 & 1.398 & 64,0\% & 0,0\% \\
BioASQ & Mistral & Fixed Top-5 & 0,733 & 0,940 & 0,349 & 0,447 & 0,761 & 2.247 & 42,2\% & 0,0\% \\
BioASQ & Mistral & Fixed Top-7 & 0,798 & 0,940 & 0,329 & 0,434 & 0,737 & 3.040 & 21,8\% & 0,0\% \\
BioASQ & Mistral & Fixed Top-10 & 0,868 & 0,940 & 0,257 & 0,346 & 0,606 & 3.889 & 0,0\% & 0,0\% \\
BioASQ & Mistral & Heuristic Rules & 0,761 & 0,940 & 0,333 & 0,445 & 0,745 & 2.560 & 34,2\% & 0,0\% \\
BioASQ & Mistral & Safe Adaptive Context & 0,825 & 0,940 & 0,339 & 0,447 & 0,758 & 2.011 & 48,3\% & 0,0\% \\
BioASQ & Llama-70B & No Retrieval & 0,000 & 0,000 & 0,197 & 0,215 & 0,508 & 120 & 96,6\% & 0,0\% \\
BioASQ & Llama-70B & Fixed Top-3 & 0,637 & 0,940 & 0,340 & 0,482 & 0,780 & 1.172 & 66,5\% & 0,0\% \\
BioASQ & Llama-70B & Fixed Top-5 & 0,733 & 0,940 & 0,354 & 0,513 & 0,782 & 1.860 & 46,8\% & 0,0\% \\
BioASQ & Llama-70B & Fixed Top-7 & 0,798 & 0,940 & 0,341 & 0,493 & 0,779 & 2.524 & 27,8\% & 0,0\% \\
BioASQ & Llama-70B & Fixed Top-10 & 0,868 & 0,940 & 0,344 & 0,492 & 0,774 & 3.496 & 0,0\% & 0,0\% \\
BioASQ & Llama-70B & Heuristic Rules & 0,761 & 0,940 & 0,349 & 0,505 & 0,780 & 2.126 & 39,2\% & 0,0\% \\
BioASQ & Llama-70B & Safe Adaptive Context & 0,825 & 0,940 & 0,342 & 0,499 & 0,779 & 1.668 & 52,3\% & 2,0\% \\
\bottomrule
\end{tabular}%
}
\end{table}


\subsection{Diagnostic Tables}
\label{app:diagnostics}

Table~\ref{tab:table_7_dataset_level_summary_of_safe_ad} summarizes average Safe Adaptive behavior across datasets.

\begin{table}[htbp]
\centering
\small
\caption{Dataset-level summary of Safe Adaptive Context. F1 retention is measured relative to Fixed Top-10 and averaged across Mistral and Llama-70B.}
\label{tab:table_7_dataset_level_summary_of_safe_ad}
\resizebox{\linewidth}{!}{%
\begin{tabular}{lrrrr}
\toprule
Dataset & Avg. Token Reduction & Avg. F1 Retention & Avg. Fallback & Interpretation \\
\midrule
SciFact & 73,3\% & 112,7\% & 5,0\% & Compact evidence is often enough \\
HotpotQA & 30,1\% & 100,4\% & 0,0\% & Multi-hop questions need broader context \\
BioASQ & 50,3\% & 115,7\% & 1,0\% & Biomedical factoids benefit from focused evidence \\
\bottomrule
\end{tabular}%
}
\end{table}

\subsubsection{Compression and Budget Diagnostic}

Table~\ref{tab:appendix_table_a4_full_compression_and_budget} compares full context, compressed context, heuristic budgeting, and Safe Adaptive Context for SciFact and HotpotQA.
\begin{table}[H]
\centering
\small
\caption{Full compression and budget diagnostic across SciFact and HotpotQA.}
\label{tab:appendix_table_a4_full_compression_and_budget}
\resizebox{\linewidth}{!}{%
\begin{tabular}{lllrrr}
\toprule
Dataset & Method & F1 & Semantic & Tokens & Token Reduction \\
\midrule
SciFact & Fixed Full Context & 0,262 & 0,569 & 3.412 & 0,0\% \\
SciFact & Compressed Fixed Full Context & 0,257 & 0,571 & 2.150 & 37,0\% \\
SciFact & Heuristic Rules & 0,266 & 0,575 & 2.234 & 34,5\% \\
SciFact & Heuristic Rules + Compact Evidence & 0,262 & 0,575 & 1.453 & 57,4\% \\
SciFact & Safe Adaptive Context & 0,265 & 0,569 & 933 & 72,6\% \\
HotpotQA & Fixed Full Context & 0,765 & 0,768 & 1.537 & 0,0\% \\
HotpotQA & Compressed Fixed Full Context & 0,752 & 0,755 & 1.430 & 6,9\% \\
HotpotQA & Heuristic Rules & 0,708 & 0,711 & 945 & 38,5\% \\
HotpotQA & Heuristic Rules + Compact Evidence & 0,696 & 0,700 & 902 & 41,3\% \\
HotpotQA & Safe Adaptive Context & 0,737 & 0,740 & 1.066 & 30,6\% \\
\bottomrule
\end{tabular}%
}
\end{table}


Token reduction is computed relative to Fixed Top-10 under the same dataset and retriever.

\subsubsection{Cross-Encoder Reranking}
Table~\ref{tab:appendix_table_a5_full_tf_idf_and_cross_encod} reports full method-level TF-IDF and cross-encoder results for Llama-70B.
\begin{table}[H]
\centering
\small
\caption{Full TF-IDF and cross-encoder reranking results for Llama-70B.}
\label{tab:appendix_table_a5_full_tf_idf_and_cross_encod}
\resizebox{\linewidth}{!}{%
\begin{tabular}{lllrrrrrrr}
\toprule
Dataset & Retriever & Method & nDCG@10 & MRR@10 & F1 & Semantic & Tokens & Token Reduction & Fallback \\
\midrule
SciFact & TF-IDF & Fixed Top-10 & 0,602 & 0,563 & 0,262 & 0,569 & 3.412 & 0,0\% & 0,0\% \\
SciFact & TF-IDF & Heuristic Rules & 0,587 & 0,557 & 0,266 & 0,575 & 2.234 & 34,5\% & 0,0\% \\
SciFact & TF-IDF & Safe Adaptive Context & 0,560 & 0,542 & 0,265 & 0,569 & 933 & 72,6\% & 5,0\% \\
SciFact & Cross-enc. & Fixed Top-10 & 0,649 & 0,629 & 0,261 & 0,573 & 3.412 & 0,0\% & 0,0\% \\
SciFact & Cross-enc. & Heuristic Rules & 0,639 & 0,626 & 0,267 & 0,569 & 1.829 & 46,4\% & 0,0\% \\
SciFact & Cross-enc. & Safe Adaptive Context & 0,624 & 0,620 & 0,265 & 0,574 & 939 & 72,5\% & 4,0\% \\
HotpotQA & TF-IDF & Fixed Top-10 & 0,742 & 0,812 & 0,765 & 0,768 & 1.537 & 0,0\% & 0,0\% \\
HotpotQA & TF-IDF & Heuristic Rules & 0,691 & 0,809 & 0,708 & 0,711 & 945 & 38,5\% & 0,0\% \\
HotpotQA & TF-IDF & Safe Adaptive Context & 0,742 & 0,812 & 0,737 & 0,740 & 1.066 & 30,6\% & 0,0\% \\
HotpotQA & Cross-enc. & Fixed Top-10 & 0,829 & 0,943 & 0,764 & 0,838 & 1.500 & 0,0\% & 0,0\% \\
HotpotQA & Cross-enc. & Heuristic Rules & 0,804 & 0,943 & 0,767 & 0,835 & 804 & 46,4\% & 0,0\% \\
HotpotQA & Cross-enc. & Safe Adaptive Context & 0,829 & 0,943 & 0,751 & 0,819 & 1.109 & 26,1\% & 0,0\% \\
BioASQ & TF-IDF & Fixed Top-10 & 0,868 & 0,940 & 0,344 & 0,774 & 3.496 & 0,0\% & 0,0\% \\
BioASQ & TF-IDF & Heuristic Rules & 0,761 & 0,940 & 0,349 & 0,780 & 2.126 & 39,2\% & 0,0\% \\
BioASQ & TF-IDF & Safe Adaptive Context & 0,825 & 0,940 & 0,342 & 0,779 & 1.668 & 52,3\% & 2,0\% \\
BioASQ & Cross-enc. & Fixed Top-10 & 0,891 & 0,960 & 0,338 & 0,773 & 3.497 & 0,0\% & 0,0\% \\
BioASQ & Cross-enc. & Heuristic Rules & 0,766 & 0,960 & 0,340 & 0,775 & 1.750 & 50,0\% & 0,0\% \\
BioASQ & Cross-enc. & Safe Adaptive Context & 0,672 & 0,960 & 0,342 & 0,776 & 814 & 76,7\% & 3,0\% \\
\bottomrule
\end{tabular}%
}
\end{table}


\subsection{Human Annotation Protocol}
\label{app:annotation-protocol}
\begin{table}[H]
\centering
\small
\caption{Human annotation label definitions.}
\label{tab:annotation_labels}
\resizebox{\linewidth}{!}{%
\begin{tabular}{ll}
\toprule
Label & Criterion \\
\midrule
CORRECT & Contains the required information and does not contradict the reference. \\
PARTIALLY\_CORRECT & Contains some required information but is incomplete, vague, or missing a condition. \\
INCORRECT & Contradicts the reference or gives a different answer. \\
NOT\_ENOUGH\_INFO & Cannot be confidently judged from the query, reference, and evidence shown. \\
\bottomrule
\end{tabular}%
}
\end{table}


\begin{table}[H]
\centering
\small
\caption{Human annotation sample construction.}
\label{tab:annotation_sample}
\begin{tabular}{llp{0.55\linewidth}}
\toprule
Dataset & Items & Selection \\
\midrule
SciFact & 40 & Stratified examples from Mistral and Llama-70B \\
HotpotQA & 40 & Examples covering the HotpotQA output distribution \\
BioASQ & 40 & Stratified examples from Mistral and Llama-70B \\
\bottomrule
\end{tabular}
\end{table}





\end{document}